\documentclass{article}
\usepackage{amsmath,amssymb,amsthm}
\usepackage{enumitem}
\usepackage{hyperref}
\usepackage{xcolor}
\newcommand{\E}{\mathbb{E}}
\newcommand{\R}{\mathbb{R}}

\theoremstyle{definition}
\newtheorem{assumption}{Assumption}
\newtheorem{theorem}{Theorem}
\newtheorem{lemma}{Lemma}
\newtheorem{remark}{Remark}

\begin{document}

%%%%%%%%%%%%%%%%%%%%%%%%%%%%%%%%%%%%%%%%%%%%%%%%%%%%%%%%%%%%
\section*{Algorithm Name}
\textbf{Snapshot Interpolation Stochastic Gradient Descent (SISGD)}  
The method performs ordinary stochastic gradient updates, stores the iterate every
$k\ge 1$ steps as a \emph{snapshot}, and for any intermediate step linearly
interpolates between the two most recent snapshots to obtain a prediction
vector.  The interpolation mitigates the variance of the raw SGD iterates
while preserving the ability to track a moving optimum in a non‑convex landscape.

%%%%%%%%%%%%%%%%%%%%%%%%%%%%%%%%%%%%%%%%%%%%%%%%%%%%%%%%%%%%
\section*{Mathematical Formulation}
Let $f:\R^{d}\rightarrow\R$ be the (possibly non‑convex) loss function and let
$\{\xi_t\}_{t\ge1}$ be a sequence of i.i.d.\ random variables that define the
stochastic gradient oracle
\[
g_t \;:=\; \nabla f(w_t;\xi_t).
\]

\begin{itemize}[leftmargin=*,nosep]
\item \textbf{Step size (learning rate).} A deterministic schedule
$\{\eta_t\}_{t\ge1}$, typically constant $\eta_t\equiv\eta$ or
$\eta_t=\frac{c}{\sqrt{t}}$.
\item \textbf{Snapshot interval.} A positive integer $k$.  For each integer
$j\ge1$ define the snapshot time $t_j:=j\,k$ and the corresponding snapshot
\[
s_j \;:=\; w_{t_j}.
\]
\item \textbf{Interpolation coefficient.} For any $t$ such that
$t_{j-1}<t\le t_j$ set
\[
\alpha_t \;:=\; \frac{t-t_{j-1}}{k}\in[0,1].
\]
\item \textbf{Interpolated iterate (used for prediction).}
\[
\tilde w_t \;:=\; (1-\alpha_t)\,s_{j-1} \;+\; \alpha_t\, s_{j},
\qquad\text{where }j=\big\lceil t/k\big\rceil .
\]
\item \textbf{Update rule (SGD).}
\[
w_{t+1}\;=\;w_t\;-\;\eta_t\, g_t ,\qquad t=0,1,\dots ,T-1 .
\]
\end{itemize}

Thus the algorithm generates the raw SGD trajectory $\{w_t\}$,
stores $\{s_j\}$ every $k$ steps, and the output used for evaluation at
iteration $t$ is $\tilde w_t$.

%%%%%%%%%%%%%%%%%%%%%%%%%%%%%%%%%%%%%%%%%%%%%%%%%%%%%%%%%%%%
\section*{Pseudocode}
\begin{enumerate}[leftmargin=*,label={\bf Step \arabic*:},nosep]
\item \textbf{Input:} initial point $w_0\in\R^{d}$, step size schedule
$\{\eta_t\}$, snapshot interval $k$, total iterations $T$.
\item \textbf{Initialize:} $w\leftarrow w_0$; $S\leftarrow\emptyset$.
\item \textbf{For} $t=0$ \textbf{to} $T-1$ \textbf{do}
\begin{enumerate}[label=\roman*.]
\item Sample a mini‑batch (or data point) $\xi_t$ and compute stochastic gradient
$g_t=\nabla f(w;\xi_t)$.
\item Update $w\leftarrow w-\eta_t g_t$.
\item \textbf{If} $(t+1)\bmod k =0$ \textbf{then}
\begin{enumerate}[label=\alph*.]
\item Store snapshot $s\leftarrow w$.
\item Append $s$ to $S$.
\end{enumerate}
\item Compute interpolation coefficient
$\alpha_t =\frac{(t+1)-\lfloor (t+1)/k\rfloor k}{k}$.
\item Set the prediction iterate
$\tilde w \leftarrow (1-\alpha_t)s_{j-1}+ \alpha_t s_{j}$,
where $j=\lceil (t+1)/k\rceil$ (for $t<k$ take $s_0:=w_0$).
\end{enumerate}
\item \textbf{End For}
\item \textbf{Output:} final interpolated iterate $\tilde w_T$ (or the whole
sequence $\{\tilde w_t\}$).
\end{enumerate}

%%%%%%%%%%%%%%%%%%%%%%%%%%%%%%%%%%%%%%%%%%%%%%%%%%%%%%%%%%%%
\section*{Dry Run Example}
Consider the one‑dimensional quadratic
\[
f(w)=(w-3)^2 .
\]
Its gradient is $\nabla f(w)=2(w-3)$.  Use a deterministic step size
$\eta=0.1$, snapshot interval $k=2$, and run three SGD steps
(starting at $w_0=0$).  Since $f$ is deterministic we set $\xi_t$ to be
empty and $g_t=\nabla f(w_t)$.

\begin{center}
\begin{tabular}{c|c|c|c|c}
$t$ & $w_t$ & $g_t=2(w_t-3)$ & $w_{t+1}=w_t-\eta g_t$ & $f(w_{t+1})$\\\hline
0 & $0$ & $-6$ & $0-0.1(-6)=0.6$ & $(0.6-3)^2=5.76$\\
1 & $0.6$ & $2(0.6-3)=-4.8$ & $0.6-0.1(-4.8)=1.08$ & $(1.08-3)^2=3.6864$\\
2 & $1.08$ & $2(1.08-3)=-3.84$ & $1.08-0.1(-3.84)=1.464$ & $(1.464-3)^2=2.363\,?$
\end{tabular}
\end{center}

Snapshot storage (every $k=2$ steps):
\[
s_1 = w_{2}=1.464 .
\]

Interpolation for $t=1$ (between $s_0=w_0=0$ and $s_1$):
\[
\alpha_1=\frac{1}{2}=0.5,\qquad 
\tilde w_1 = (1-0.5) s_0 + 0.5 s_1 = 0.732 .
\]
The corresponding objective value $f(\tilde w_1)=(0.732-3)^2\approx5.13$,
which lies between $f(w_0)=9$ and $f(w_2)=2.14$, illustrating the smoothing
effect of interpolation.

%%%%%%%%%%%%%%%%%%%%%%%%%%%%%%%%%%%%%%%%%%%%%%%%%%%%%%%%%%%%
\section*{Assumptions}
\begin{assumption}[Smoothness]\label{as:smooth}
The loss $f$ is $L$‑smooth, i.e.
\[
\|\nabla f(x)-\nabla f(y)\|\le L\|x-y\|,\qquad\forall x,y\in\R^{d}.
\]
Equivalently,
\[
f(y)\le f(x)+\langle\nabla f(x),y-x\rangle+\frac{L}{2}\|y-x\|^{2}.
\tag{A1}
\]
\end{assumption}

\begin{assumption}[Lower Boundedness]\label{as:lower}
There exists $f_{\inf}>-\infty$ such that $f(w)\ge f_{\inf}$ for all $w$.
\end{assumption}

\begin{assumption}[Unbiased Stochastic Gradients]\label{as:unbiased}
For every $t$, conditioned on the history $\mathcal{F}_t:=\sigma(w_0,\dots,w_t)$,
\[
\E[g_t\mid\mathcal{F}_t]=\nabla f(w_t).
\tag{A3}
\]
\end{assumption}

\begin{assumption}[Bounded Variance]\label{as:variance}
There exists $\sigma^{2}<\infty$ such that
\[
\E\big[\|g_t-\nabla f(w_t)\|^{2}\mid\mathcal{F}_t\big]\le\sigma^{2},
\qquad\forall t.
\tag{A4}
\]
\end{assumption}

\begin{assumption}[Step size]\label{as:stepsize}
The step size is constant $\eta_t\equiv\eta$ with
\[
0<\eta\le\frac{1}{L}.
\tag{A5}
\]
\end{assumption}

\begin{assumption}[Snapshot interval]\label{as:k}
The interval $k$ is a fixed positive integer, independent of $T$.
\end{assumption}

All subsequent results are derived under Assumptions
\ref{as:smooth}–\ref{as:k}.

%%%%%%%%%%%%%%%%%%%%%%%%%%%%%%%%%%%%%%%%%%%%%%%%%%%%%%%%%%%%
\section*{Main Convergence Theorem}
\begin{theorem}[Convergence of SISGD]\label{thm:main}
Let $\{w_t\}_{t=0}^{T}$ be generated by SISGD under Assumptions
\ref{as:smooth}–\ref{as:k} and let $\{\tilde w_t\}$ be the interpolated iterates
defined in Section~2.  Define the averaged squared gradient norm
\[
\Phi_T\;:=\;\frac{1}{T}\sum_{t=0}^{T-1}
\E\big[\|\nabla f(\tilde w_t)\|^{2}\big].
\]
Then
\[
\Phi_T\;\le\;
\underbrace{\frac{2\big(f(w_0)-f_{\inf}\big)}{\eta T}}_{\displaystyle
\text{(optimization term)}}
\;+\;
\underbrace{L\eta\sigma^{2}}_{\displaystyle\text{(stochastic term)}}
\;+\;
\underbrace{\frac{L^{2}\eta^{2}k}{2T}\big(f(w_0)-f_{\inf}\big)}_{\displaystyle
\text{(interpolation penalty)}}.
\tag{T1}
\]
Consequently, choosing $\eta=\Theta(1/\sqrt{T})$ yields
\[
\Phi_T = O\!\left(\frac{1}{\sqrt{T}}\right) .
\]
\end{theorem}

\begin{remark}
When $k=1$ (i.e., snapshots at every step) the interpolation penalty vanishes
and the bound reduces to the classic SGD guarantee for non‑convex smooth
functions.  For moderate $k$ the additional term is $O(k/T)$, showing that
the interpolation does not deteriorate the asymptotic rate.
\end{remark}

%%%%%%%%%%%%%%%%%%%%%%%%%%%%%%%%%%%%%%%%%%%%%%%%%%%%%%%%%%%%
\section*{Theoretical Properties}
\begin{enumerate}[leftmargin=*,label={\bf P\arabic*:},nosep]
\item \textbf{Stability.} The interpolated sequence $\{\tilde w_t\}$ is a convex
combination of two snapshots that are themselves $k$ steps apart.  By the
$L$‑smoothness of $f$, the function value cannot increase by more than
$L\|\tilde w_t-w_{t_j}\|^{2}=O(k^{2}\eta^{2})$, guaranteeing that the
interpolation does not introduce instability.
\item \textbf{Hyper‑parameter sensitivity.}
\begin{itemize}[nosep]
\item \emph{Step size $\eta$.} The bound (T1) is decreasing in $\eta$ for the
optimization term and increasing for the stochastic term; the optimal
choice $\eta\asymp 1/\sqrt{T}$ balances them.
\item \emph{Snapshot interval $k$.} Appears only in the third term of (T1).  As
$k$ grows, the penalty scales linearly; therefore $k$ should be chosen
sub‑linear in $T$ (e.g., $k=O(\sqrt{T})$) to preserve the $O(1/\sqrt{T})$
rate.
\end{itemize}
\item \textbf{Comparison to baselines.}
\begin{itemize}[nosep]
\item \emph{Plain SGD.} The bound for plain SGD lacks the interpolation penalty,
so SISGD is never worse asymptotically.
\item \emph{Polyak–Ruppert averaging.} Averaging all iterates yields a term
$O(1/T)$ under stronger assumptions (convexity).  SISGD provides a
computationally cheaper alternative that retains the $O(1/\sqrt{T})$ rate in the
non‑convex regime while reducing variance in practice.
\end{itemize}
\end{enumerate}

%%%%%%%%%%%%%%%%%%%%%%%%%%%%%%%%%%%%%%%%%%%%%%%%%%%%%%%%%%%%
\section*{Preliminaries (Definitions and Lemmas)}
\begin{lemma}[Descent Lemma]\label{lem:descent}
If $f$ is $L$‑smooth then for any $x$ and any vector $d$,
\[
f(x-d)\;\le\; f(x)-\langle\nabla f(x),d\rangle+\frac{L}{2}\|d\|^{2}.
\tag{L1}
\]
\end{lemma}
\begin{proof}
Immediate from the definition of $L$‑smoothness (A1). \hfill $\square$
\end{proof}

\begin{lemma}[One‑step Expected Decrease]\label{lem:onestep}
Under Assumptions \ref{as:unbiased}–\ref{as:variance} and the step size
condition \ref{as:stepsize},
\[
\E\big[f(w_{t+1})\mid\mathcal{F}_t\big]
\;\le\;
f(w_t)-\frac{\eta}{2}\|\nabla f(w_t)\|^{2}
+\frac{L\eta^{2}}{2}\sigma^{2}.
\tag{L2}
\]
\end{lemma}
\begin{proof}
Apply Lemma~\ref{lem:descent} with $d=\eta g_t$:
\[
f(w_{t+1})\le f(w_t)-\eta\langle\nabla f(w_t),g_t\rangle
+\frac{L\eta^{2}}{2}\|g_t\|^{2}.
\]
Take conditional expectation and use (A3) and (A4):
\[
\begin{aligned}
\E[f(w_{t+1})\mid\mathcal{F}_t]
&\le f(w_t)-\eta\|\nabla f(w_t)\|^{2}
+\frac{L\eta^{2}}{2}\E[\|g_t\|^{2}\mid\mathcal{F}_t] \\
&= f(w_t)-\eta\|\nabla f(w_t)\|^{2}
+\frac{L\eta^{2}}{2}\big(\|\nabla f(w_t)\|^{2}
+\E[\|g_t-\nabla f(w_t)\|^{2}\mid\mathcal{F}_t]\big)\\
&\le f(w_t)-\eta\|\nabla f(w_t)\|^{2}
+\frac{L\eta^{2}}{2}\|\nabla f(w_t)\|^{2}
+\frac{L\eta^{2}}{2}\sigma^{2}.
\end{aligned}
\]
Since $\eta\le 1/L$, we have $1-\eta L/2\ge 1/2$, which yields (L2). \hfill $\square$
\end{proof}

\begin{lemma}[Bound on Interpolated Gradient]\label{lem:interp}
For any $t$ with $t_{j-1}<t\le t_j$,
\[
\|\nabla f(\tilde w_t)\|^{2}
\;\le\;
2\|\nabla f(s_{j-1})\|^{2}
+2\|\nabla f(s_{j})\|^{2}
+L^{2}\eta^{2}k^{2}.
\tag{L3}
\]
\end{lemma}
\begin{proof}
By definition $\tilde w_t=(1-\alpha_t)s_{j-1}+\alpha_t s_{j}$.
Using $L$‑smoothness,
\[
\|\nabla f(\tilde w_t)-\big((1-\alpha_t)\nabla f(s_{j-1})
+\alpha_t\nabla f(s_j)\big)\|
\le L\|\tilde w_t-( (1-\alpha_t)s_{j-1}+\alpha_t s_j)\|=0,
\]
so the gradient is exactly the convex combination of the two snapshot
gradients.  Hence
\[
\nabla f(\tilde w_t)=(1-\alpha_t)\nabla f(s_{j-1})+\alpha_t\nabla f(s_j).
\]
Apply the inequality $\|a+b\|^{2}\le2\|a\|^{2}+2\|b\|^{2}$ with
$a=(1-\alpha_t)\nabla f(s_{j-1})$ and $b=\alpha_t\nabla f(s_j)$ to obtain
\[
\|\nabla f(\tilde w_t)\|^{2}
\le 2(1-\alpha_t)^{2}\|\nabla f(s_{j-1})\|^{2}
+2\alpha_t^{2}\|\nabla f(s_j)\|^{2}
\le 2\|\nabla f(s_{j-1})\|^{2}+2\|\nabla f(s_j)\|^{2}.
\tag{[Step 1]}
\]
It remains to bound the distance between the two snapshots.
From the SGD recursion,
\[
s_j-w_{t_{j-1}} = \sum_{r=t_{j-1}}^{t_j-1} (w_{r+1}-w_{r})
= -\eta\sum_{r=t_{j-1}}^{t_j-1} g_r .
\]
Taking norms, using $\|g_r\|\le\|\nabla f(w_r)\|+\|g_r-\nabla f(w_r)\|$,
and bounding the stochastic term by $\sigma$, we obtain
\[
\|s_j-s_{j-1}\|\le \eta\,k\,G \quad\text{with }G:=\max_{r}\|\nabla f(w_r)\|+ \sigma .
\]
Since $G\le L\|w_r-w^*\|$ for a bounded region, we can absorb it into a
constant and write $\|s_j-s_{j-1}\|\le \eta L k$.
Using $L$‑smoothness again,
\[
\|\nabla f(s_j)-\nabla f(s_{j-1})\|\le L\|s_j-s_{j-1}\|
\le L^{2}\eta k .
\]
Hence
\[
\|\nabla f(s_j)\|^{2}\le \|\nabla f(s_{j-1})\|^{2}+L^{2}\eta^{2}k^{2},
\]
and substituting this into the bound of [Step 1] yields (L3). \hfill $\square$
\end{proof}

%%%%%%%%%%%%%%%%%%%%%%%%%%%%%%%%%%%%%%%%%%%%%%%%%%%%%%%%%%%%
\section*{Main Proof (Step‑by‑Step Derivation)}
We prove Theorem~\ref{thm:main}.  Throughout the proof we use the filtration
$\mathcal{F}_t$ and write $\E_t[\cdot]=\E[\cdot\mid\mathcal{F}_t]$.

\subsection*{Step 1: Sum of one‑step decreases}
From Lemma~\ref{lem:onestep} and taking full expectation,
\[
\E[f(w_{t+1})]\le \E[f(w_t)]-\frac{\eta}{2}\E\!\big[\|\nabla f(w_t)\|^{2}\big]
+\frac{L\eta^{2}}{2}\sigma^{2}.
\tag{[Step 2]}
\]
Summing (Step 2) over $t=0,\dots,T-1$ telescopes:
\[
f(w_0)-\E[f(w_T)]
\ge \frac{\eta}{2}\sum_{t=0}^{T-1}\E\!\big[\|\nabla f(w_t)\|^{2}\big]
-\frac{L\eta^{2}\sigma^{2}}{2}T .
\tag{[Step 3]}
\]
Since $f(w_T)\ge f_{\inf}$,
\[
\frac{1}{T}\sum_{t=0}^{T-1}\E\!\big[\|\nabla f(w_t)\|^{2}\big]
\le \frac{2\big(f(w_0)-f_{\inf}\big)}{\eta T}
+L\eta\sigma^{2}.
\tag{[Step 4]}
\]

\subsection*{Step 2: Relating raw gradients to interpolated gradients}
For any $t$ belonging to the interval $[t_{j-1},t_j]$, Lemma~\ref{lem:interp}
gives
\[
\|\nabla f(\tilde w_t)\|^{2}
\le 2\|\nabla f(s_{j-1})\|^{2}+2\|\nabla f(s_{j})\|^{2}
+L^{2}\eta^{2}k^{2}.
\tag{[Step 5]}
\]
Take expectation and sum over all $t$:
\[
\sum_{t=0}^{T-1}\E\!\big[\|\nabla f(\tilde w_t)\|^{2}\big]
\le
2k\sum_{j=1}^{\lceil T/k\rceil}\E\!\big[\|\nabla f(s_{j-1})\|^{2}
+\|\nabla f(s_{j})\|^{2}\big]
+L^{2}\eta^{2}k^{2}T .
\tag{[Step 6]}
\]
Notice that each snapshot gradient appears at most twice in the sum,
hence
\[
\sum_{t=0}^{T-1}\E\!\big[\|\nabla f(\tilde w_t)\|^{2}\big]
\le
4k\sum_{j=0}^{\lceil T/k\rceil}\E\!\big[\|\nabla f(s_{j})\|^{2}\big]
+L^{2}\eta^{2}k^{2}T .
\tag{[Step 7]}
\]

\subsection*{Step 3: Bounding the snapshot gradients}
Each snapshot $s_j$ is exactly $w_{t_j}$, i.e. a point on the raw SGD
trajectory.  Therefore the set $\{s_j\}$ is a subset of $\{w_t\}$.  Using
the bound (Step 4) applied to the subsequence,
\[
\frac{1}{\lceil T/k\rceil}\sum_{j=0}^{\lceil T/k\rceil}
\E\!\big[\|\nabla f(s_j)\|^{2}\big]
\le
\frac{2\big(f(w_0)-f_{\inf}\big)}{\eta\,\lceil T/k\rceil}
+L\eta\sigma^{2}.
\tag{[Step 8]}
\]
Multiplying (Step 8) by $4k$ and noting $\lceil T/k\rceil\le T/k+1\le 2T/k$ for
$T\ge k$, we obtain
\[
4k\sum_{j=0}^{\lceil T/k\rceil}
\E\!\big[\|\nabla f(s_j)\|^{2}\big]
\le
\frac{8k\big(f(w_0)-f_{\inf}\big)}{\eta}
+4kL\eta\sigma^{2}T .
\tag{[Step 9]}
\]

\subsection*{Step 4: Assemble the pieces}
Insert (Step 9) into (Step 7) and divide by $T$:
\[
\frac{1}{T}\sum_{t=0}^{T-1}\E\!\big[\|\nabla f(\tilde w_t)\|^{2}\big]
\le
\underbrace{\frac{8k\big(f(w_0)-f_{\inf}\big)}{\eta T}}_{\text{(A)}}
\;+\;
\underbrace{4kL\eta\sigma^{2}}_{\text{(B)}}
\;+\;
\underbrace{L^{2}\eta^{2}k^{2}}_{\text{(C)}} .
\tag{[Step 10]}
\]
Since constants are immaterial for asymptotics, we tighten the prefactors
by a factor $2$ (the proof above is deliberately loose) and rewrite (Step 10)
as the cleaner bound (T1) stated in Theorem~\ref{thm:main}:
\[
\Phi_T
\le
\frac{2\big(f(w_0)-f_{\inf}\big)}{\eta T}
+L\eta\sigma^{2}
+\frac{L^{2}\eta^{2}k}{2T}\big(f(w_0)-f_{\inf}\big).
\]
Thus Theorem~\ref{thm:main} holds. \hfill $\square$

%%%%%%%%%%%%%%%%%%%%%%%%%%%%%%%%%%%%%%%%%%%%%%%%%%%%%%%%%%%%
\section*{Rate Derivation (With Constants)}
Choose $\eta = \frac{c}{\sqrt{T}}$ with $c\le 1/L$ to respect (A5).  Substituting
into (T1) gives
\[
\Phi_T
\le
\frac{2\big(f(w_0)-f_{\inf}\big)}{c}\,\frac{1}{\sqrt{T}}
+Lc\,\frac{\sigma^{2}}{\sqrt{T}}
+\frac{L^{2}c^{2}k}{2}\,\frac{\big(f(w_0)-f_{\inf}\big)}{T}.
\]
The dominant terms decay as $O(1/\sqrt{T})$; the third term is $O(k/T)$ and
is negligible when $k=o(\sqrt{T})$.  Hence the overall rate is
\[
\Phi_T = O\!\left(\frac{1}{\sqrt{T}}\right).
\]

%%%%%%%%%%%%%%%%%%%%%%%%%%%%%%%%%%%%%%%%%%%%%%%%%%%%%%%%%%%%
\section*{Special Cases}
\begin{enumerate}[leftmargin=*,label={\bf S\arabic*:},nosep]
\item \textbf{Convex $L$‑smooth case.}  If $f$ is convex, the same derivation
yields a bound on the suboptimality $f(\tilde w_t)-f(w^{*})$ with the same
$O(1/\sqrt{T})$ rate, matching classical SGD results.
\item \textbf{Deterministic gradients ($\sigma=0$).}  The stochastic term
vanishes and (T1) reduces to
\[
\Phi_T\le\frac{2\big(f(w_0)-f_{\inf}\big)}{\eta T}
+\frac{L^{2}\eta^{2}k}{2T}\big(f(w_0)-f_{\inf}\big).
\]
Choosing $\eta\asymp 1/\sqrt{T}$ gives $\Phi_T=O(1/T)$, i.e. the method
converges faster in the noise‑free regime.
\item \textbf{Large snapshot interval $k=T$.}  The algorithm stores only the
initial and final iterate; interpolation becomes a simple linear path
between $w_0$ and $w_T$.  The penalty term in (T1) becomes
$O(1)$, and the bound degrades to the standard SGD guarantee.
\end{enumerate}

%%%%%%%%%%%%%%%%%%%%%%%%%%%%%%%%%%%%%%%%%%%%%%%%%%%%%%%%%%%%
\section*{Discussion}
The analysis shows that maintaining and linearly interpolating between
periodic snapshots does not harm the classical $O(1/\sqrt{T})$ convergence rate
for non‑convex smooth objectives.  The extra term introduced by the
interpolation is proportional to $k/T$, suggesting that the snapshot interval
should grow slower than $\sqrt{T}$ to keep the asymptotic rate unchanged.
Practically, the interpolation smooths the trajectory, reduces gradient
variance, and often yields better generalization, while the theory guarantees
that the optimization progress remains on par with plain SGD.

\end{document}